% Level 2, Buffer B2 -- Generative Art Jam (can be used any week after week 3)
\kcpart{B2}{Generative Art Jam}{Rules plus randomness make art that can surprise even its own artist.}
\lessoninfo{Concepts: lists (R), random (R), functions with parameters and return (R), random seed (I, light), generative art (I)}{Tools: Thonny, dice, squared paper}{CSTA: 2-AP-12, 2-AP-14, 2-AP-16}

\begin{teachernote}[When to use this session]
Any time after week 3 (lists): after a missed week, before a holiday, or as a playful break. It
also gives Art Studio project teams a head start. Nothing new and hard: lists, \texttt{random},
functions and loops, used for fun.
\end{teachernote}

\begin{objectives}
  \item The student can use a list of colors and random choices to make a picture that changes every run. (L2.G4, L2.G7)
  \item The student can write a drawing function with parameters and call it many times from a loop. (L2.G3)
  \item The student can use \texttt{random.seed} to get the same \enquote{random} picture again. (L2.G1, L2.G7)
  \item The student can give credit when remixing someone else's art idea or code. (L2.G9)
\end{objectives}

\begin{materials}
  \item Print: nothing required. Optional: this page's challenge cards, cut out (or write them on index cards).
  \item One die per pair, squared paper, colored pencils (or plain pencils and the pattern key below).
  \item On each Mac: Thonny, \menu{File > New}, save as \texttt{art\_jam.py} in \texttt{Documents}.
  \item Prep (10 minutes): type the guided demo program once yourself and run it.
\end{materials}

\needspace{16\baselineskip}
\section{At a glance}
\begin{timeplan}
  \timeslot{0--8}{Warm-up}{Dice art on squared paper: the same rules make different pictures.}
  \timeslot{8--18}{Concept}{Generative art = rules + randomness. A seed lets you get the same picture back.}
  \timeslot{18--28}{Guided demo}{Type the random flower garden: a palette list, a flower function, a list of recipes, a seed.}
  \timeslot{28--52}{Hands-on build}{Art jam with challenge cards; artist and engineer roles swap every 8 minutes; bug hunt at minute 36.}
  \timeslot{52--60}{Wrap-up}{Gallery: run each program twice (same seed, new seed). Credits, reflection, exit question.}
\end{timeplan}

\section{Warm-up (0--8 min)}
\begin{unplugged}[Dice art]
Each student draws a 6 × 6 grid on squared paper. For every square, roll the die once and follow
the rule key (write it on the board):
\textbf{1} = leave empty, \textbf{2} = dot, \textbf{3} = diagonal line /, \textbf{4} = diagonal line \textbackslash,
\textbf{5} = filled square, \textbf{6} = circle. Color optional: pick from orange, sky blue, green, blue.
After 5 minutes, compare: everyone used the same rules, yet every picture is different.
Then try this: one student reads out a list of rolls (\enquote{3, 3, 5, 1…}) and everyone follows the
same list. Now the pictures match! That list of rolls is what programmers call a \textbf{seed}.
With 1 student, you both draw and compare.
\end{unplugged}

\section{Concept (8--18 min)}
\begin{teacherscript}
“Generative art is art made by rules plus randomness. You write the rules: draw a flower, use these
colors, put it somewhere on the screen. The computer rolls the dice: where exactly, how big, which
color. Every run is a brand-new picture.

We need three things you already know. A list of colors, our palette. A function with parameters,
like \texttt{flower(t, x, y, size, color)}, so we can draw a flower anywhere, any size. And a loop to
draw lots of them. One new trick: \texttt{random.seed(7)}. It is like reading out the same dice rolls:
the random numbers come out the same every time, so if you love a picture, write down its seed and
you can get it back.”
\end{teacherscript}

\begin{analogy}[A recipe with a handful of sprinkles]
A cake recipe says \enquote{throw on a handful of sprinkles}. Same recipe, but every cake looks a
bit different, because the sprinkles land in new places. If you filmed exactly where each sprinkle
landed, you could make the very same cake again. The recipe is your code, the sprinkles are
\texttt{random}, and the film is the seed.
\end{analogy}

\section{Guided demo (18--28 min)}
\begin{workedexample}[A random flower garden]
Type the program below in this order, running it (\keys{F5}) after steps 2 and 4:
\begin{enumerate}
  \item The \texttt{PALETTE} list (the Okabe--Ito colors we use all year).
  \item \texttt{flower}: penup, go to (x, y), then one petal per loop (two arcs). Test it once with \texttt{flower(pen, 0, 0, 100, "\#E69F00")}.
  \item \texttt{petal\_count}: a function that returns a number, so bigger flowers get more petals.
  \item \texttt{make\_garden}: it does not draw anything. It returns a list of flower recipes (x, y, size, color). Say: \enquote{Plan first, draw later: the list is the plan.}
  \item The main part: draw every recipe in the list. \texttt{screen.tracer(0)} and \texttt{screen.update()} make it appear in one go.
\end{enumerate}
Run it twice: the same garden appears, because of \texttt{seed=7}. Change it to \texttt{seed=8}: a new garden.
\end{workedexample}

\begin{pycode}[art\_jam.py -- random flower garden (tested)]
# Generative Art Jam - starter idea
import random
import turtle

PALETTE = ["#E69F00", "#56B4E9", "#009E73", "#F0E442",
           "#0072B2", "#D55E00", "#CC79A7"]

def petal_count(size):
    """Bigger flowers get more petals (between 5 and 12)."""
    return max(5, min(12, size // 10))

def flower(t, x, y, size, color):
    """Draw one flower made of loops, centered at (x, y)."""
    t.penup()
    t.goto(x, y)
    t.pendown()
    t.color(color)
    petals = petal_count(size)
    for _ in range(petals):
        t.circle(size / 2, 60)
        t.left(120)
        t.circle(size / 2, 60)
        t.left(120 + 360 / petals)   # finish the petal, then turn to the next one

def make_garden(how_many, seed):
    """Return a list of flower recipes: (x, y, size, color)."""
    random.seed(seed)  # same seed = same picture every time
    recipes = []
    for _ in range(how_many):
        x = random.randint(-300, 300)
        y = random.randint(-200, 200)
        size = random.randint(30, 120)
        recipes.append((x, y, size, random.choice(PALETTE)))
    return recipes

if __name__ == "__main__":
    screen = turtle.Screen()
    screen.bgcolor("black")
    screen.tracer(0)          # draw everything at once
    pen = turtle.Turtle()
    pen.hideturtle()
    pen.pensize(2)
    for x, y, size, color in make_garden(25, seed=7):
        flower(pen, x, y, size, color)
    screen.update()
    turtle.done()
\end{pycode}

\section{Hands-on build (28--52 min)}
\begin{pairroles}[Artist and engineer]
The \textbf{artist} decides what the picture should look like and describes it in words
(\enquote{mirror it}, \enquote{tiny flowers at the top}). The \textbf{engineer} types and runs it.
Swap every 8 minutes. \textbf{1 student:} you are the artist and the student engineers; swap once.
\textbf{3:} add a \textbf{seed keeper} who writes down every seed and what it looked like. \textbf{4:} two pairs.
\end{pairroles}

\textbf{Challenge cards} (pick one, finish it, pick another):
\begin{enumerate}
  \item \textbf{Mirror world:} draw every flower twice, at \texttt{(x, y)} and at \texttt{(-x, y)}.
  \item \textbf{Mood palette:} make a second list with only three colors and use it instead.
  \item \textbf{Growing garden:} in the main loop, make flowers lower on the screen bigger (use y to choose the size).
  \item \textbf{New shape:} write \texttt{star(t, x, y, size, color)} (5 points: forward, right 144) and mix stars with flowers.
  \item \textbf{Seed hunt:} try ten seeds and write down your three favorites.
\end{enumerate}

\begin{bughunt}[Three broken gardens (minute 36)]
Write each bug on the board. Pairs predict what will happen, then fix it.
\begin{enumerate}
  \item \texttt{random.seed(seed)} moved \emph{inside} the \texttt{for} loop. All 25 flowers land in exactly the same spot (one flower!). Fix: seed once, before the loop.
  \item \texttt{x = random.randint(300, -300)}: \texttt{ValueError: empty range\ldots} (the exact words depend on the Python version). Fix: smaller number first.
  \item \texttt{PALETTE = ["E69F00", …]} (no \texttt{\#}): \texttt{TurtleGraphicsError: bad color string: E69F00}. Fix: color codes start with \texttt{\#}.
\end{enumerate}
\end{bughunt}

\begin{freebuild}[Jam time (last 10 minutes of the build)]
Free choice: combine challenge cards, invent a new shape function, or make art for a project title
screen. The only rule: keep the palette list and a function with parameters.
\end{freebuild}

\section{Wrap-up (52--60 min)}
Gallery: each pair runs its program with the same seed twice, then with a new seed. Everyone names
one thing they like in someone else's art.

\begin{formativecheck}[Exit question]
\textbf{Look for:} a palette list used with \texttt{random.choice}; a drawing function called with
different arguments from a loop.
\textbf{Ask:} \enquote{What does \texttt{random.seed(7)} do?}
\textbf{Answer:} it makes the random numbers come out the same every run, so you get the same picture
again. A different seed gives a different picture.
\end{formativecheck}

\begin{reflection}
Which rule made the biggest difference to your picture? What surprised you? If you could print one of
today's pictures as a poster, which seed would you choose?
\end{reflection}

\begin{safety}[Credit for generative art]
Generative artists share ideas all the time, and they give credit. If you remix a friend's idea or
code, add a comment such as \texttt{\# Mirror idea from Mei}, using a nickname. Seeds are a safe thing to
share: they are only numbers and contain no personal information. Screenshots of your art stay on the
family computer unless a parent decides otherwise.
\end{safety}

\section{Notes for the teacher}
\begin{commonmistakes}
  \item \textbf{The drawing is very slow.} \texttt{screen.tracer(0)} is missing, or \texttt{screen.update()} is not called at the end (then nothing appears).
  \item \textbf{Nothing changes between runs.} The seed is fixed: that is correct! Change the number, or use \texttt{seed=None} (Python then seeds from the clock) to get a new garden every run.
  \item \textbf{Mixing up parameters}: \texttt{flower(pen, size, x, y, color)} puts flowers in odd places. Read the \texttt{def} line aloud and match each argument.
\end{commonmistakes}

\begin{differentiation}
\supportpath Use the demo program unchanged and only change the numbers: \texttt{how\_many}, the seed,
the size range, the background. Then try challenge card 2.
\extensionpath Save favorite seeds to a file \texttt{favorites.txt} (one per line, with
\texttt{"a"} mode) and write a menu that redraws a saved favorite: a mini Art Studio.
\end{differentiation}

\begin{stretchchallenge}[Rules that react]
Make each flower's color depend on where it is: flowers on the left use the first half of the palette,
flowers on the right the second half. Use \texttt{if x < 0:} and list slices \texttt{PALETTE[:3]} and
\texttt{PALETTE[3:]}.
\end{stretchchallenge}

\begin{teachernote}[If you are short on time]
Skip the dice art and show the analogy instead. Drop \texttt{petal\_count} from the demo (use a fixed
\texttt{petals = 8}). Keep at least 10 minutes of jam time and the 2-minute gallery: that is the fun part.
\end{teachernote}
